\documentclass[10pt,a4paper]{amsart}
\usepackage[margin=19mm]{geometry}
\usepackage{amsmath,amssymb,mathtools}
\usepackage[scr=rsfs,cal=euler]{mathalfa}
\usepackage{xfrac}
\usepackage[unicode,hidelinks,bookmarks=false]{hyperref}
\newcommand{\ie}{i.e.\ }
\newcommand{\cf}{cf.\ }
\newcommand{\resp}{respectively\ }
\newcommand{\Subsection}{\subsection}
\providecommand{\nobreakdash}{\nobreak\hbox{-}}
\setlength{\emergencystretch}{2em}
\multlinegap0pt
\usepackage{bm}
\usepackage{bbm}
\usepackage{dsfont}
\usepackage{xspace}
\usepackage{cancel}
\usepackage{mathtools}
\usepackage{amsthm}
\makeatletter
\@ifundefined{theorem}{\newtheorem{theorem}{Theorem}}{}
\@ifundefined{proposition}{\newtheorem{proposition}[theorem]{Proposition}}{}
\makeatother
\DeclareMathOperator{\codis}{codis}
\DeclareMathOperator{\dist}{dist}
\DeclareMathOperator{\dis}{dis}
\DeclareMathOperator{\dom}{dom}
\DeclareMathOperator{\ecc}{ecc}
\DeclareMathOperator{\gr}{Gr}
\DeclareMathOperator{\rad}{rad}
\DeclareMathOperator{\sep}{sep}
\newcommand{\GH}{d_{GH}}
\newcommand{\MP}[1]{(A_{#1},\chi_{#1})}
\newcommand{\N}{\mathbb{N}}
\newcommand{\R}{\mathbb{R}}
\newcommand{\Tau}{\mathcal T}
\newcommand{\setdef}[2]{\left\{{#1} \ \middle|\ {#2}\right\}}
\title[Gromov-Hausdorff distance between chromatic metric pairs and]{Gromov-Hausdorff distance between chromatic metric pairs and stability of the six-pack}
\author{Ond\v{r}ej Draganov, Sophie Rosenmeier, Nicol\`o Zava}
\date{Source: arXiv:2507.17994v1 (CC BY 4.0)}
\begin{document}
\maketitle

\noindent\emph{Source and licence.} Gromov-Hausdorff distance between chromatic metric pairs and stability of the six-pack. Version arXiv:2507.17994v1 (CC BY 4.0), distributed under the Creative Commons Attribution 4.0 International licence, as recorded in the arXiv licence metadata for this exact version and shown on the abstract page https://arxiv.org/abs/2507.17994v1. Full rights notice: licenses/arxiv-2507.17994-CC-BY-4.0.txt.\par\smallskip

\noindent\textit{Editorial scope.} Counted unit: the Proposition whose source label is \texttt{prop:chromatic\_invariants} in the article (source0.tex lines 1154--1204, source label \texttt{prop:chromatic\_invariants}), delivered together with the article's own complete original proof of it (source0.tex lines 1205--1234, one \texttt{proof} environment of 30 lines). No mathematics is written, repaired or re-derived: statement and proof are the article's own text line for line. Required context retained uncounted: none. Every definition, display and label of those lines is retained unchanged. Omitted from this package and not redistributed: 1--1153, 1235--2315. Nothing inside a retained excerpt is omitted or altered: every retained body is delivered exactly as the article's lines are written, so no omission receipt of any kind accompanies this package. Citations: the retained excerpts carry no citation command at all, so no bibliography entry of the article is transcribed and no reference is redistributed. Reference adaptation: none was needed and none was made; every reference inside the delivered unit resolves inside it, so no unresolved reference or citation is produced. Printed slips kept literal: no printed word, symbol, hypothesis or number of the counted statement or of its proof was altered. Layout and class: the article's own class is \texttt{article} with packages including inputenc, indentfirst, amsmath, amsfonts, amssymb, amsthm, amscd, ifthen, hyperref, latexsym, mathrsfs, algorithm, algorithmic, paralist; the wrapper is the lane's pinned \texttt{amsart} editorial wrapper at 10pt on A4, which restates the theorem-like environments the retained text needs and adds the article's own macro definitions that the retained text needs (\texttt{\textbackslash GH}, \texttt{\textbackslash MP}, \texttt{\textbackslash N}, \texttt{\textbackslash R}, \texttt{\textbackslash Tau}, \texttt{\textbackslash codis}, \texttt{\textbackslash dis}, \texttt{\textbackslash dist}, \texttt{\textbackslash dom}, \texttt{\textbackslash ecc}, \texttt{\textbackslash gr}, \texttt{\textbackslash rad}, \texttt{\textbackslash sep}, \texttt{\textbackslash setdef}), with extra wrapper packages (bm, bbm, dsfont, xspace, cancel, mathtools). The delivered numbering is the unit's own. Assets: no retained span contains a figure environment, a graphics-inclusion command, a PDF-page inclusion, a table or a TikZ picture; no image file is copied into this package, the package declares no asset dependency and no image file is redistributed with it. Licence: version arXiv:2507.17994v1 of this article is distributed under the Creative Commons Attribution 4.0 International licence, as recorded in the arXiv licence metadata for this exact version and shown on the abstract page https://arxiv.org/abs/2507.17994v1. A full rights notice is delivered at licenses/arxiv-2507.17994-CC-BY-4.0.txt. Characters: every retained excerpt is pure ASCII, so the delivered engine, XeLaTeX with its default Latin Modern text font, reports no missing character.
\begin{proposition}\label{prop:chromatic_invariants}
    Let $\N\in C\subseteq\mathcal P(\N)$. For every pair of chromatic metric pairs $\MP{1}$ and $\MP{2}$, and every pair of distinct $\sigma,\tau\in\mathcal \Tau_C$, we consider the following invariant
    \[
        d_{\mathcal{L}}^{\sigma,\tau}((A_1,\chi_1), (A_2, \chi_2))
        =
        \inf_{R\in\mathcal R(\chi_1^{-1}(\sigma),\chi_2^{-1}(\sigma))}\sup_{(x_1,x_2)\in R}d_H(\mathcal L_{\MP{1}}^{\sigma,\tau}(x_1),\mathcal L_{\MP{2}}^{\sigma,\tau}(x_2)).
    \]
    Then, the following inequalities hold:
    \begin{equation}
    \begin{aligned}
        \lvert\rad_{\MP{1}}^{\sigma,\tau}-\rad_{\MP{2}}^{\sigma,\tau}\rvert&\,
        \leq
        d_H(\mathcal E_{\MP{1}}^{\sigma,\tau},\mathcal E_{\MP{2}}^{\sigma,\tau})
        = \\
        %
        &\,=
        \inf_{R\in\mathcal R(\chi_1^{-1}(\sigma),\chi_2^{-1}(\sigma))}\sup_{(x_1,x_2)\in R}\lvert\ecc_{\MP{1}}^{\sigma,\tau}(x_1)-\ecc_{\MP{2}}^{\sigma,\tau}(x_2)\rvert\leq \\
        %
        &\,\leq
        d_{\mathcal{L}}^{\sigma,\tau}((A_1,\chi_1), (A_2, \chi_2)),
    \end{aligned}
    \end{equation}
    \begin{equation}
    \begin{aligned}
        \lvert\dist_{\MP{1}}^{\sigma,\tau}-\dist_{\MP{2}}^{\sigma,\tau}\rvert&\,
        \leq
        d_H(\mathcal S_{\MP{1}}^{\sigma,\tau},\mathcal S_{\MP{2}}^{\sigma,\tau})
        =\\
        %
        &\,=
        \inf_{R\in\mathcal R(\chi_1^{-1}(\sigma),\chi_2^{-1}(\sigma))}\sup_{(x_1,x_2)\in R}\lvert\sep_{\MP{1}}^{\sigma,\tau}(x_1)-\sep_{\MP{2}}^{\sigma,\tau}(x_2)\rvert
        \leq \\
        %
        &\,\leq
        d_{\mathcal{L}}^{\sigma,\tau}((A_1,\chi_1), (A_2, \chi_2)),
    \end{aligned}
    \end{equation}
    
    \begin{equation}\label{eq:stable_chromatic_invariants}
    \begin{aligned}
    d_H(\mathcal D_{\MP{1}}^{\sigma,\tau},\mathcal D_{\MP{2}}^{\sigma,\tau})
    &\,\leq
    d_{\mathcal{L}}^{\sigma,\tau}((A_1,\chi_1), (A_2, \chi_2))
    \leq \\
    %
    &\,\leq
    2\GH^C(\chi_1,\chi_2)\leq 2\GH^C(\MP{1},\MP{2}),
    \end{aligned}
    \end{equation}
    where, we recall, $\chi_i$ denotes the chromatic metric space $(\dom\chi_i, \chi_i)$.
\end{proposition}

\begin{proof}
Most of the inequalities are easy to verify---relying on the observation that the difference between the infima (or suprema) of two subsets of $\R$ are always upper-bounded by their Hausdorff distance. For the equalities, a correspondence on either side can be used to construct a correspondence on the other such that the differences considered under the suprema are the same.

Let us focus on the most relevant inequality for our purposes, i.e., the second one in \eqref{eq:stable_chromatic_invariants}. Fix $\varepsilon>0$, and two $C$-constrained maps $f\colon\chi_1\to\chi_2$ and $g\colon\chi_2\to\chi_1$ 
such that
\[
    \max\{\dis f,\dis g,\codis(f,g)\}\leq 2\GH^C(\chi_1,\chi_2)+\varepsilon.
\]
Define $R=\gr(f|_{\chi_1^{-1}(\sigma)})\cup(\gr(g|_{\chi_2^{-1}(\sigma)}))^{-1}$, which is a correspondence between $\chi_1^{-1}(\sigma)$ and $\chi_2^{-1}(\sigma)$.
Fix $(x_1,x_2)\in R$. By definition, either $f(x_1)=x_2$ or $x_1=g(x_2)$. Let us consider the former case; the latter can be treated analogously. We upper-bound the Hausdorff distance between $\mathcal L_{\MP{1}}^{\sigma,\tau}(x_1)$ and $\mathcal L_{\MP{2}}^{\sigma,\tau}(x_2)$ by considering the correspondence \[
    S = \setdef{(d_1(x_1, x'), d_2(x_2, f(x')))}{x'\in\chi_1^{-1}(\tau)} \cup \setdef{(d_1(x_1, g(y')), d_2(x_2, y'))}{y'\in\chi_2^{-1}(\tau)}.
\] 
For every $x'\in\chi_1^{-1}(\tau)$ and $y'\in\chi_2^{-1}(\tau)$, we have
\begin{align*}
&\lvert d_1(x_1,x')-d_2(x_2,f(x'))\rvert = \lvert d_1(x_1,x')-d_2(f(x_1),f(x'))\rvert \leq \dis f,
\\
&\lvert d_1(x_1,g(y'))-d_2(x_2,y')\rvert = \lvert d_1(x_1,g(y'))-d_2(f(x_1),y')\rvert \leq \codis(f,g),
\end{align*}
and so 
\begin{equation}\label{eq:stability}
    d_H(\mathcal L_{\MP{1}}^{\sigma,\tau}(x_1),\mathcal L_{\MP{2}}^{\sigma,\tau}(x_2))\leq\max\{\dis f,\codis(f,g)\}.
\end{equation}
For the analogous inequality in the case where $x_1=g(x_2)$, we need to use $\dis g$ in the argument of the maximum on the right-hand side of \eqref{eq:stability}.

Hence, we have shown that
\[
    \sup_{(x_1,x_2)\in R}d_H(\mathcal L_{\MP{1}}^{\sigma,\tau}(x_1),\mathcal L_{\MP{2}}^{\sigma,\tau}(x_2))\leq\max\{\dis f,\dis g,\codis(f,g)\}\leq 2\GH^C(\chi_1,\chi_2)+\varepsilon.
\]
Since $\varepsilon$ can be taken arbitrarily small, we have obtained the desired inequality.
\end{proof}

\end{document}
